% 03_ostrich.tex - where to get OSTRICH, the version used, and the loader workaround
\section{Get OSTRICH}

\step{Download OSTRICH and test it}\label{st:ostrich}
OSTRICH (by L.~Shawn Matott) is not part of this workshop. Get it from its official sources:
\begin{itemize}
\item The University of Waterloo: Linux, macOS and Windows programs, and the Raven--OSTRICH tutorial files with
  Demo~C1:\newline \url{https://www.civil.uwaterloo.ca/envmodelling/Ostrich.html}
\item The author's page: the source code, which you can compile with \cmd{make GCC}:\newline
  \url{https://www.eng.buffalo.edu/~lsmatott/Ostrich/OstrichMain.html}
\item The documentation of the current versions: \url{https://usbr.github.io/ostrich}
\end{itemize}
This workshop was made with \cmd{Ostrich\_v20171219\_GCC\_Linux.exe} from the folder \cmd{Demo\_C1} of the Raven--OSTRICH
tutorial files. It reports itself as \emph{OSTRICH version \V{t6:ostver}}. Unzip the Raven--OSTRICH tutorial files in a
folder of their own and go into that folder. Then copy the program into a folder \cmd{\textasciitilde/ostrich}:
\begin{shell}
$ mkdir -p ~/ostrich
$ cp Demo_C1/Ostrich_v20171219_GCC_Linux.exe ~/ostrich/
$ chmod +x ~/ostrich/Ostrich_v20171219_GCC_Linux.exe
\end{shell}
This program was built on a computing cluster. It asks for a program loader (the part of Linux that starts programs) that
other Linux systems do not have. Start it through your system's loader instead. Put the whole command in the variable
\cmd{OSTRICH}, and test it in the folder \cmd{\textasciitilde/ostrich}, which has no OSTRICH input file:
\begin{shell}
$ export OSTRICH="/lib64/ld-linux-x86-64.so.2 $HOME/ostrich/Ostrich_v20171219_GCC_Linux.exe"
$ cd ~/ostrich
$ $OSTRICH
Starting up MPI required 0.000002 seconds
InitDataLine() : Couldn't open file |ostIn.txt|
Ostrich Error Report for Processor 0
A total of 1 errors and/or warnings were reported
FILE I/O ERROR             : GetInFileName()(): couldn't open |ostIn.txt|
...
\end{shell}
\see these lines (the time on the first line differs). OSTRICH started and complained only that \cmd{ostIn.txt} is missing.
That is correct here.
\begin{trouble}
\begin{itemize}
\item \cmd{cannot execute: required file not found}, although the file is there: you started the program directly. Use the
  \cmd{export OSTRICH=...} line above, which starts it through the loader. This is how it ran for this workshop.
\item \cmd{/lib64/ld-linux-x86-64.so.2: No such file or directory}: your Linux keeps the loader elsewhere. Find it with
  \cmd{ls /lib*/ld-linux*} and use that path. Or compile OSTRICH from its source code.
\end{itemize}
\end{trouble}
